\documentclass[11pt,a4paper]{article}
\usepackage[utf8]{inputenc}
\usepackage[T1]{fontenc}
\usepackage[brazil]{babel}
\usepackage{lmodern}
\usepackage[margin=2cm]{geometry}
\usepackage{amsmath,amssymb}
\usepackage{booktabs}
\usepackage{graphicx}
\usepackage{hyperref}
\usepackage{xcolor}
\usepackage{titlesec}
\usepackage{tabularx}
\usepackage{enumitem}
\usepackage{float}
\usepackage{array}
\renewcommand{\tabularxcolumn}[1]{>{\raggedright\arraybackslash}p{#1}}
\usepackage{newunicodechar}
\usepackage[most]{tcolorbox}
\usepackage{etoolbox}

\hypersetup{colorlinks=true, linkcolor=blue!70!black, citecolor=blue!70!black, urlcolor=blue!70!black}

\titleformat{\section}{\large\bfseries}{\thesection.}{0.5em}{}
\titleformat{\subsection}{\normalsize\bfseries}{\thesubsection.}{0.5em}{}
\titlespacing*{\section}{0pt}{1.2ex plus 0.5ex minus .2ex}{0.6ex}
\titlespacing*{\subsection}{0pt}{1.0ex plus 0.4ex minus .2ex}{0.4ex}
\setlength{\parskip}{0.45em}
\setlength{\parindent}{0pt}
\setlist{itemsep=0.15em,topsep=0.3em}

% --- símbolos ---
\newunicodechar{π}{\ensuremath{\pi}}
\newunicodechar{τ}{\ensuremath{\tau}}
\newunicodechar{φ}{\ensuremath{\varphi}}
\newunicodechar{ψ}{\ensuremath{\psi}}
\newunicodechar{Σ}{\ensuremath{\Sigma}}
\newunicodechar{→}{\ensuremath{\rightarrow}}
\newunicodechar{⇒}{\ensuremath{\Rightarrow}}
\newunicodechar{∈}{\ensuremath{\in}}
\newunicodechar{−}{\ensuremath{-}}
\newunicodechar{≈}{\ensuremath{\approx}}
\newunicodechar{≠}{\ensuremath{\neq}}
\newunicodechar{≡}{\ensuremath{\equiv}}
\newunicodechar{≤}{\ensuremath{\leq}}
\newunicodechar{≥}{\ensuremath{\geq}}
\newunicodechar{′}{\ensuremath{'}}
\newunicodechar{‑}{-}
\newunicodechar{✔}{\ensuremath{\checkmark}}
\newunicodechar{✘}{\ensuremath{\times}}

% --- cores e caixas ---
\definecolor{cC}{HTML}{7A4FA3}
\definecolor{cS}{HTML}{2F7D74}
\definecolor{cAcc}{HTML}{2F5D8A}
\newcommand{\metrictag}[2]{{\setlength{\fboxsep}{1.5pt}%
  \ifstrequal{#1}{c}{\colorbox{cC}{\textcolor{white}{\scriptsize\textbf{#2}}}}{%
  \ifstrequal{#1}{s}{\colorbox{cS}{\textcolor{white}{\scriptsize\textbf{#2}}}}{%
  \colorbox{black!60}{\textcolor{white}{\scriptsize\textbf{#2}}}}}}}
\newcommand{\pagechip}[1]{{\setlength{\fboxsep}{1.5pt}\colorbox{cAcc}{\textcolor{white}{\scriptsize\texttt{#1}}}}}
\newcommand{\readtag}[1]{\texttt{\small #1}}
\newcommand{\partbar}[1]{\par\vspace{1.2em}\noindent\colorbox{black!80}{\parbox{\dimexpr\linewidth-2\fboxsep}{\color{white}\bfseries\sffamily\MakeUppercase{#1}}}\par\vspace{0.4em}}
\newtcolorbox{citacard}[1]{breakable, enhanced, colback=white, colframe=black!25, boxrule=0.5pt, arc=2pt,
  left=6pt, right=6pt, top=4pt, bottom=4pt, title={\small #1}, coltitle=black, colbacktitle=black!6, fonttitle=\normalfont}
\newenvironment{origquote}{\par\itshape\small}{\par}
\newenvironment{tradquote}{\par\small\textsc{Tradução:}\ }{\par}
\newcommand{\explab}[1]{\par\smallskip{\small\color{cAcc}\textbf{\MakeUppercase{#1}}}\par}
\colorlet{notebg}{blue!4}\colorlet{notefg}{cAcc}
\colorlet{okbg}{green!5}\colorlet{okfg}{green!45!black}
\colorlet{warnbg}{orange!8}\colorlet{warnfg}{orange!80!black}
\newtcolorbox{infobox}[2][note]{breakable, enhanced, boxrule=0pt, arc=2pt, left=6pt, right=6pt, top=4pt, bottom=4pt,
  colback=#1bg, borderline west={3pt}{0pt}{#1fg}, frame hidden,
  before upper={\ifstrempty{#2}{}{\textbf{#2}\par}}}
\colorlet{pairc}{cC}\colorlet{pairs}{cS}
\newenvironment{paircol}[2]{\begin{tcolorbox}[width=0.485\linewidth, nobeforeafter, box align=top, boxrule=0.5pt,
  colframe=black!25, colback=white, borderline north={3pt}{0pt}{pair#1}, left=5pt, right=5pt,
  before upper={{\color{pair#1}\small\textbf{\MakeUppercase{#2}}}\par}]}{\end{tcolorbox}}

\newunicodechar{ℓ}{\ensuremath{\ell}}

\begin{document}

\begin{center}
    {\LARGE \textbf{Verificação da referência do GradientSHAP}} \\[0.35em]
    {\large ECGs normais, referência balanceada e 200 sorteios} \\[0.3em]
    \textbf{Código:} \texttt{src/baseline\_check.py}, \texttt{src/baseline\_balanced.py} \\[0.2em]
    \textbf{Resultados:} \texttt{results/modeloB\_basecheck/}, \texttt{results/modeloB\_basebal/} \\[0.2em]
    \textbf{Aluno:} Arthur Santos de Oliveira (RA: 156297) \\[0.2em]
    \textbf{Atualizado em:} 26/10/2026 (primeira versão: 25/10/2026) \\[0.2em]
    \textbf{Repositório:} \url{https://github.com/Arthur-so/shap-ecg-arrhythmia}
\end{center}
\vspace{-0.5em}
\hrule
\vspace{0.8em}

\textbf{Por que este documento.} Ao estudar a complexidade e o Random Logit, surgiram duas dúvidas sobre a forma como os mapas SHAP do TCC são gerados: (1) os ECGs de referência usados pelo GradientSHAP (as \textbf{baselines}) são uma mistura de arritmias sorteada do treino, quando clinicamente a referência natural seria o ECG \textbf{normal}; (2) cada mapa é a média de só \textbf{50 sorteios}, e o erro dessa média (o \textbf{piso}) é grande. Este documento descreve dois testes que medem o peso das duas escolhas e termina com as opções de decisão. O primeiro comparou a referência atual com uma referência de ECGs normais; o segundo, feito depois, testou uma referência \textbf{balanceada por classe} e a referência atual com 200 sorteios, para separar o efeito da referência do efeito dos sorteios.

\textbf{Como ler.} A \textbf{Parte I} explica o que é a referência e por que ela importa; a \textbf{Parte II} descreve o primeiro teste; a \textbf{Parte III} traz os seus resultados; a \textbf{Parte IV} traz o segundo teste e compara as quatro opções; a \textbf{Parte V}, as opções para a decisão. Todos os números são médias sobre os 622 mapas do modelo B (predições corretas do teste) e usam as 36 regiões do batimento (12 derivações × pré-QRS, QRS e ST-T), como as demais métricas do TCC.

\begin{infobox}[note]{Resumo}
\begin{itemize}
\item \textbf{200 sorteios cortam o erro pela metade.} O piso (distância entre dois mapas que só diferem nos sorteios) cai de 0,33 para 0,17, como prevê a estatística de uma média.
\item \textbf{A referência normal muda os mapas, mas moderadamente.} Fora SNR, o mapa com referência normal e 200 sorteios tem cosseno médio de 0,92 com o mapa atual. Trocar só o sorteio das 32 baselines por outro aleatório já muda o mapa quase tanto.
\item \textbf{Com a referência normal, a classe SNR perde o sentido:} explicar um ECG normal em relação a ECGs normais deixa quase nada a explicar (soma média do mapa $-1{,}5$, contra 16,2 hoje).
\item \textbf{Prévia com os mapas novos:} a Consistency sobe na maioria das classes (global 0,42 $\to$ 0,54); a complexidade praticamente não muda.
\item \textbf{Só aumentar os sorteios quase não muda as métricas} (referência atual com 200 sorteios: Consistency 0,42 $\to$ 0,44; complexidade igual). A subida da Consistency vem principalmente da troca de referência, e não da redução do ruído.
\item \textbf{A referência balanceada} (4 ECGs de cada classe) fica entre a atual e a normal: mantém SNR interpretável, corrige o peso excessivo de BRD e muda os mapas além do ruído, mas pouco (Consistency 0,46).
\item \textbf{Nada foi trocado.} Os resultados oficiais continuam com os mapas atuais. A decisão de adotar uma nova referência é do aluno com o orientador (Parte V).
\end{itemize}
\end{infobox}

\partbar{Parte I · A referência do GradientSHAP}

\section{O que é a referência}
O GradientSHAP não explica a saída do modelo “em absoluto”. Ele explica a \textbf{diferença} entre a saída para o ECG e a saída média para um conjunto de ECGs de referência, as \textbf{baselines}. A soma de todas as importâncias do mapa é, aproximadamente,
\[
\text{soma do mapa} \approx f_c(\text{ECG}) - \text{média de } f_c(\text{baselines}),
\]
em que $f_c$ é o logit da classe explicada $c$. Cada importância vem de termos do tipo gradiente × (ECG $-$ baseline). Por isso, uma região do ECG que é \textbf{parecida} com a mesma região das baselines recebe pouca importância, mesmo que o modelo a use: o fator (ECG $-$ baseline) é pequeno ali.

\begin{infobox}[note]{Consequência}
O mapa responde à pergunta: \textit{o que neste ECG difere da referência e leva o modelo à classe $c$?} Trocar a referência troca a pergunta.
\end{infobox}

\section{As 32 baselines atuais}
Em \texttt{src/explain.py}, a referência são 32 ECGs sorteados do conjunto de treino (4\,812 ECGs), sempre com a semente 42 do projeto. O número 32 é o padrão do código, sem justificativa formal registrada. O sorteio ficou desbalanceado por acaso:
\begin{center}\small
\begin{tabular}{lrrrrrrrrr}
\toprule
& \textbf{SNR} & \textbf{AF} & \textbf{IAVB} & \textbf{BRE} & \textbf{BRD} & \textbf{CAP} & \textbf{CVP} & \textbf{STD} & \textbf{STE} \\
\midrule
\% no treino & 13,3 & 17,8 & 10,5 & 3,4 & 27,1 & 9,0 & 10,2 & 12,7 & 3,2 \\
Esperado em 32 & 4,3 & 5,7 & 3,3 & 1,1 & 8,7 & 2,9 & 3,3 & 4,0 & 1,0 \\
\textbf{Nas 32 atuais} & 2 & 4 & 5 & \textbf{0} & \textbf{14} & 2 & 5 & 1 & \textbf{0} \\
\bottomrule
\end{tabular}
\end{center}
Nenhuma baseline é individualmente extrema (duração de 10 a 30 s, amplitude máxima abaixo do percentil 99 do treino), mas BRD está super-representado e BRE e STE estão ausentes.

\section{Por que o ECG normal}
\begin{center}\small
\begin{tabularx}{\linewidth}{lX}
\toprule
\textbf{Referência} & \textbf{Pergunta que o mapa responde} \\
\midrule
Atual: mistura de ECGs do treino & O que neste ECG faz o modelo dar à classe uma saída diferente da que ele dá, em média, aos ECGs do treino (a maioria com alguma arritmia)? \\ \addlinespace
ECGs normais (SNR) & O que neste ECG, em relação a um coração saudável, levou o modelo à classe? \\
\bottomrule
\end{tabularx}
\end{center}
A primeira é o padrão do SHAP (uma amostra dos dados como referência) e não está errada, mas compara a arritmia com outras arritmias. A segunda é mais próxima do raciocínio clínico: uma arritmia se define pelo que a afasta do ritmo normal. A escolha da referência é uma questão conhecida na literatura de explicabilidade (ver, por exemplo, Sturmfels, Lundberg e Lee, 2020, \textit{Visualizing the Impact of Feature Attribution Baselines}, Distill; não fichado).

\textbf{Um efeito colateral previsto.} Com a referência normal, a classe SNR perde o sentido: explicar um ECG normal em relação a ECGs normais deixa a soma do mapa perto de zero, e o mapa vira quase só ruído.

\section{Os 50 sorteios e o piso}
O GradientSHAP estima o mapa como a média de \texttt{n\_samples} $= 50$ termos sorteados (qual baseline, um ponto no caminho entre ela e o ECG e um pequeno ruído). Com outros 50 sorteios, sai um mapa um pouco diferente. O \textbf{piso} é a distância entre dois mapas do mesmo ECG e da mesma classe que só diferem nos sorteios; na sensibilidade, foi 0,355. O erro de uma média cai aproximadamente com $1/\sqrt{n}$: com 200 sorteios (4 vezes mais), o piso deveria cair pela metade. A explicação detalhada está na seção 6.2 do fichamento da complexidade (\texttt{fichamentos/bhatt-2020-complexity}).

\partbar{Parte II · O teste}

\section{Desenho}
Para cada um dos 622 mapas (ECG × classe explicada), o GradientSHAP foi refeito em sete versões, todas com a \textbf{mesma semente} de sorteios internos, para que as diferenças venham só da referência ou do número de sorteios:
\begin{center}\small
\begin{tabularx}{\linewidth}{lXl}
\toprule
\textbf{Mapa} & \textbf{Referência (32 ECGs do treino)} & \textbf{Sorteios} \\
\midrule
atual & a de \texttt{explain.py} (semente 42) & 50, semente $s$ \\
atual$'$ & a de \texttt{explain.py} & 50, semente $s+1$ \\
normal & 32 ECGs com rótulo único SNR (semente 42) & 50, semente $s$ \\
normal$'$ & a mesma normal & 50, semente $s+1$ \\
aleatória & outro sorteio do treino (semente 43): SNR 6, AF 2, IAVB 6, BRE 1, BRD 7, CAP 3, CVP 2, STD 5, STE 1 & 50, semente $s$ \\
normal200 & a mesma normal; 4 blocos de 50, o 1.º é o mapa normal & 200 \\
normal200$'$ & a mesma normal; 4 blocos a partir de normal$'$ & 200 \\
\bottomrule
\end{tabularx}
\end{center}
Os 200 sorteios são feitos como média de 4 execuções de 50 com sementes diferentes, o que equivale estatisticamente a uma média de 200 e cabe na memória da GPU. As 32 normais são ECGs do \textbf{treino} (nenhum do teste) com rótulo único de SNR; todos os 642 SNR do treino têm rótulo único.

\textbf{Comparações.} Distância $\ell_2$ entre os mapas normalizados para norma 1, nas 36 regiões (0 = iguais; 2 = opostos; o cosseno entre os mapas é $1 - d^2/2$):
\begin{center}\small
\begin{tabularx}{\linewidth}{lX}
\toprule
\textbf{Comparação} & \textbf{O que mede} \\
\midrule
atual × atual$'$ & Piso atual (50 sorteios) \\
atual × aleatória & Quanto muda o mapa só por trocar o sorteio das 32 baselines (controle) \\
atual × normal & Efeito da referência normal, com 50 sorteios \\
atual × normal200 & Efeito da referência normal, com o mapa normal mais preciso \\
normal × normal$'$ & Piso com a referência normal, 50 sorteios \\
normal200 × normal200$'$ & Piso com a referência normal, 200 sorteios \\
\bottomrule
\end{tabularx}
\end{center}
Também foram registradas a soma de cada mapa e a região mais importante. Os mapas normal200 foram salvos para recalcular as métricas sem gerar os mapas de novo.

\textbf{Execução.} Kaggle (GPU T4), versão 21 do runner, 23/10/2026, 72 min ($11$ execuções de 50 sorteios por mapa, cerca de 6\,800 no total). Nenhum erro.

\section{Como ler o resultado}
\begin{itemize}
\item \textbf{atual × normal perto do piso:} a referência pesa pouco; basta justificar a escolha no texto.
\item \textbf{atual × normal bem acima do piso e acima de atual × aleatória:} a referência normal muda a explicação de fato.
\item \textbf{atual × normal $\approx$ atual × aleatória, ambos acima do piso:} os mapas dependem de qual conjunto se usa, mas não especificamente de ele ser normal.
\end{itemize}

\partbar{Parte III · Resultados}

\section{Os sorteios: 200 cortam o piso pela metade}
\begin{center}\small
\begin{tabular}{lrrrrrrrrrr}
\toprule
& \textbf{SNR} & \textbf{AF} & \textbf{IAVB} & \textbf{BRE} & \textbf{BRD} & \textbf{CAP} & \textbf{CVP} & \textbf{STD} & \textbf{STE} & \textbf{Global} \\
\midrule
Piso atual (50) & 0,41 & 0,35 & 0,36 & 0,23 & 0,32 & 0,46 & 0,32 & 0,40 & 0,35 & 0,355 \\
Piso normal (50) & 0,51 & 0,31 & 0,30 & 0,30 & 0,25 & 0,41 & 0,29 & 0,39 & 0,36 & 0,331 \\
\textbf{Piso normal (200)} & 0,26 & 0,17 & 0,15 & 0,15 & 0,12 & 0,21 & 0,15 & 0,20 & 0,15 & \textbf{0,167} \\
\bottomrule
\end{tabular}
\end{center}
Com a mesma referência normal, passar de 50 para 200 sorteios reduz o piso de 0,331 para 0,167, quase exatamente a metade prevista ($1/\sqrt{4}$), e em todas as classes. A ordem dos pisos entre as classes se mantém (CAP e STD mais altos; BRD mais baixo), o que confirma que o piso depende de como o modelo responde a cada classe, e não só do número de sorteios.

\section{A referência: muda os mapas, mas moderadamente}
\begin{center}\small
\setlength{\tabcolsep}{4pt}
\begin{tabular}{lrrrrrr}
\toprule
& \textbf{Piso} & \textbf{Atual ×} & \textbf{Atual ×} & \textbf{Atual ×} & \textbf{Cosseno} & \textbf{Acima} \\
\textbf{Classe} & \textbf{atual} & \textbf{aleatória} & \textbf{normal (50)} & \textbf{normal (200)} & \textbf{(200)} & \textbf{do piso} \\
\midrule
SNR  & 0,41 & 0,52 & 0,99 & 1,01 & 0,47 & 99\% \\
AF   & 0,35 & 0,37 & 0,41 & 0,37 & 0,92 & 54\% \\
IAVB & 0,36 & 0,36 & 0,40 & 0,37 & 0,92 & 56\% \\
BRE  & 0,23 & 0,30 & 0,37 & 0,32 & 0,95 & 75\% \\
BRD  & 0,32 & 0,36 & 0,41 & 0,39 & 0,92 & 81\% \\
CAP  & 0,46 & 0,52 & 0,54 & 0,49 & 0,87 & 67\% \\
CVP  & 0,32 & 0,36 & 0,37 & 0,34 & 0,94 & 64\% \\
STD  & 0,40 & 0,46 & 0,50 & 0,46 & 0,89 & 72\% \\
STE  & 0,35 & 0,41 & 0,44 & 0,42 & 0,91 & 86\% \\
\midrule
Sem SNR & 0,35 & 0,39 & 0,43 & 0,39 & 0,92 & 69\% \\
\bottomrule
\end{tabular}
\end{center}
“× normal (200)” é a distância entre o mapa atual e o mapa com referência normal e 200 sorteios; “Acima do piso” é a fração dos mapas em que essa distância passa do piso atual do próprio mapa.

\begin{itemize}
\item \textbf{A referência normal muda os mapas além do ruído}, em 69\% dos mapas (fora SNR). Mas a mudança é moderada: o cosseno médio entre o mapa atual e o novo é 0,92 (1 seria igual). As regiões que importam continuam, em grande parte, as mesmas.
\item \textbf{Trocar só o sorteio das 32 baselines já muda o mapa quase tanto} (0,39 contra 0,39 com a referência normal e 200 sorteios, fora SNR). Ou seja, os mapas atuais dependem de qual conjunto de referência foi sorteado; a referência normal não muda os mapas muito mais do que outro sorteio qualquer.
\item \textbf{Parte da distância é ruído do próprio mapa atual.} Mesmo um mapa normal perfeito ficaria a cerca de 0,25 do mapa atual só pelo erro dos 50 sorteios do mapa atual (piso $\div \sqrt{2}$). O efeito da referência em si é, portanto, menor do que os 0,39 sugerem.
\end{itemize}

\subsection{Região mais importante}
Fração dos mapas em que a região mais importante é a mesma no mapa atual e no mapa normal (200): 66\% no global; 92\% em BRD, 76\% em CVP, 65 a 67\% em AF, IAVB e BRE, 57\% em STE, 45 a 47\% em STD e CAP, e 29\% em SNR. Alguns deslocamentos:
\begin{center}\small
\begin{tabularx}{\linewidth}{lXX}
\toprule
\textbf{Classe} & \textbf{Referência atual} & \textbf{Referência normal (200)} \\
\midrule
BRD  & V1 QRS (85\%) & V1 QRS (91\%) \\
IAVB & aVR pré-QRS (62\%), II pré-QRS (13\%) & aVR pré-QRS (52\%), II pré-QRS (33\%) \\
STD  & aVR ST-T (37\%), II ST-T (11\%) & II ST-T (26\%), aVR ST-T (17\%) \\
STE  & aVR pré-QRS (57\%), V1 QRS (21\%) & aVR pré-QRS (29\%), V2 pré-QRS (21\%) \\
\bottomrule
\end{tabularx}
\end{center}
Esses deslocamentos são registrados como observação; a plausibilidade clínica das regiões não é avaliada no TCC.

\subsection{Soma dos mapas}
\begin{center}\small
\begin{tabular}{lrrrrrrrrr}
\toprule
& \textbf{SNR} & \textbf{AF} & \textbf{IAVB} & \textbf{BRE} & \textbf{BRD} & \textbf{CAP} & \textbf{CVP} & \textbf{STD} & \textbf{STE} \\
\midrule
Atual & 16,2 & 15,9 & 16,9 & 21,4 & 11,2 & 12,9 & 13,5 & 13,3 & 12,9 \\
Normal (200) & \textbf{$-1{,}5$} & 16,5 & 17,6 & 20,3 & \textbf{18,6} & 11,9 & 15,8 & 14,3 & 12,3 \\
\bottomrule
\end{tabular}
\end{center}
\begin{itemize}
\item \textbf{SNR:} a soma vai a $-1{,}5$, ou seja, quase nada a explicar, como previsto. Com a referência normal, os mapas de SNR não são interpretáveis (a distância ao mapa atual, perto de 1, e a baixa coincidência da região principal confirmam isso).
\item \textbf{BRD:} a soma sobe de 11,2 para 18,6. É o efeito esperado do desbalanceamento: 14 das 32 baselines atuais têm BRD, então a saída média de BRD na referência atual é alta e sobra menos diferença para explicar. Com a referência normal, a explicação de BRD passa a cobrir toda a evidência de BRD.
\item Nas demais classes, a soma muda pouco.
\end{itemize}

\section{Prévia das métricas com os mapas novos}
Os mapas normal200 permitem recalcular, sem GPU, as métricas que não chamam o modelo. Nesta prévia mudam \textbf{duas coisas ao mesmo tempo} (referência e número de sorteios), então os efeitos não podem ser separados.
\begin{center}\small
\begin{tabular}{lrrrrrrrrrr}
\toprule
& \textbf{SNR} & \textbf{AF} & \textbf{IAVB} & \textbf{BRE} & \textbf{BRD} & \textbf{CAP} & \textbf{CVP} & \textbf{STD} & \textbf{STE} & \textbf{Global} \\
\midrule
Consistency atual & 0,21 & 0,42 & 0,45 & 0,25 & 0,64 & 0,14 & 0,26 & 0,51 & 0,03 & 0,42 \\
Consistency nova & (0,58) & 0,47 & 0,59 & 0,35 & 0,76 & 0,16 & 0,50 & 0,45 & 0,03 & 0,54 \\ \addlinespace
Complexidade atual & 26,9 & 22,8 & 23,3 & 25,8 & 22,4 & 25,7 & 24,7 & 26,9 & 25,5 & 24,2 \\
Complexidade nova & (25,7) & 22,8 & 22,6 & 26,2 & 23,1 & 25,7 & 24,8 & 27,4 & 25,6 & 24,3 \\
\bottomrule
\end{tabular}
\end{center}
Consistency com $k = 2$ (valor observado); complexidade em número efetivo de regiões (de 36). Os valores de SNR estão entre parênteses porque, com a referência normal, não são interpretáveis.
\begin{itemize}
\item \textbf{Consistency:} sobe na maioria das classes (CVP 0,26 $\to$ 0,50; IAVB 0,45 $\to$ 0,59; BRD 0,64 $\to$ 0,76), cai em STD (0,51 $\to$ 0,45) e não muda em CAP e STE, que continuam as menos consistentes. É coerente com o que se esperava: com menos ruído, o par das 2 regiões principais se repete mais entre ECGs da mesma classe.
\item \textbf{Complexidade:} praticamente igual (24,2 $\to$ 24,3 no global); as classes mais concentradas continuam AF, BRD e IAVB, e a mais espalhada, STD.
\item As métricas que chamam o modelo (Faithfulness Correlation e Estimate, Max/Avg-Sensitivity, Random Logit) não foram recalculadas.
\end{itemize}

\section{Limitações do teste}
\begin{itemize}
\item Só uma referência normal (um sorteio de 32 SNR) foi testada; outro sorteio de normais também mudaria os mapas, como mostrou o controle aleatório.
\item Este primeiro teste não gerou o mapa com a referência \textbf{atual} e 200 sorteios, que separaria o efeito dos sorteios do efeito da referência nas métricas. O segundo teste (Parte IV) resolve isso.
\item A hipótese levantada no Random Logit (o mapa de BRE parece com o de muitas classes porque BRE não está nas baselines) não pode ser testada sem refazer o Random Logit com a nova referência.
\end{itemize}

\partbar{Parte IV · Segundo teste: referência balanceada}

\section{Desenho do segundo teste}
O primeiro teste deixou duas perguntas: (1) quanto da mudança veio da referência e quanto dos 200 sorteios? (2) existe uma referência que corrija o desbalanceamento da atual sem tirar o sentido de SNR? Para responder, gerou-se, para os mesmos 622 mapas e com as mesmas sementes internas:
\begin{center}\small
\begin{tabularx}{\linewidth}{lXl}
\toprule
\textbf{Mapa} & \textbf{Referência} & \textbf{Sorteios} \\
\midrule
atual200 & a atual (32 do treino) & 200 \\
balanceada200 & \textbf{4 ECGs de cada uma das 9 classes} (36, rótulo único, treino, semente 42) & 200 \\
balanceada200$'$ & a mesma, outra semente (para o piso) & 200 \\
\bottomrule
\end{tabularx}
\end{center}
Com os mapas normal200 do primeiro teste, ficam quatro opções comparáveis, todas sobre os mesmos ECGs: \textbf{A} (atual, 50 sorteios, os mapas de hoje), \textbf{B} (normal, 200), \textbf{C} (atual, 200) e \textbf{D} (balanceada, 200). Execução: Kaggle (GPU T4), versão 22 do runner, 25--26/10/2026, cerca de 85 min. Nenhum erro. As distâncias entre os mapas de 200 sorteios e as métricas foram calculadas localmente com os mapas salvos.

\section{O que muda nos mapas}
Distância média entre os mapas (36 regiões, norma 1). Como referência de ruído, dois mapas com a \textbf{mesma} referência e 200 sorteios ficam a cerca de 0,17 a 0,18 um do outro (piso com 200 sorteios: 0,167 na normal e 0,183 na balanceada).
\begin{center}\small
\setlength{\tabcolsep}{4pt}
\begin{tabular}{lrrrrrrrrrr}
\toprule
& \textbf{SNR} & \textbf{AF} & \textbf{IAVB} & \textbf{BRE} & \textbf{BRD} & \textbf{CAP} & \textbf{CVP} & \textbf{STD} & \textbf{STE} & \textbf{Sem SNR} \\
\midrule
A × C (só sorteios) & 0,25 & 0,22 & 0,21 & 0,14 & 0,20 & 0,28 & 0,20 & 0,24 & 0,22 & 0,21 \\
C × D (atual × balanceada) & 0,37 & 0,28 & 0,26 & 0,25 & 0,32 & 0,39 & 0,26 & 0,36 & 0,28 & 0,31 \\
C × B (atual × normal) & 0,99 & 0,30 & 0,31 & 0,28 & 0,34 & 0,40 & 0,27 & 0,39 & 0,36 & 0,33 \\
D × B (balanceada × normal) & 0,92 & 0,28 & 0,28 & 0,20 & 0,23 & 0,37 & 0,25 & 0,35 & 0,33 & 0,27 \\ \addlinespace
Piso da balanceada (200) & 0,23 & 0,19 & 0,18 & 0,13 & 0,13 & 0,25 & 0,17 & 0,23 & 0,18 & 0,18 \\
\bottomrule
\end{tabular}
\end{center}
\begin{itemize}
\item \textbf{Só os sorteios (A × C):} 0,21, o que se espera do erro de 50 sorteios do mapa atual. Aumentar os sorteios não muda o que o mapa diz; só o torna mais preciso.
\item \textbf{A referência muda os mapas além do ruído, mas pouco:} com 200 sorteios dos dois lados, trocar a referência atual pela balanceada dá 0,31 e pela normal, 0,33 (fora SNR), contra um ruído de cerca de 0,18.
\item \textbf{A balanceada fica entre a atual e a normal}, um pouco mais perto da normal (0,27 de distância) do que a atual está (0,33). Em BRD a aproximação é clara (0,23): a referência balanceada tira o excesso de BRD da atual.
\item \textbf{SNR continua interpretável com a balanceada} (0,37 da atual, perto das demais classes), ao contrário da normal (0,99).
\end{itemize}

\textbf{Soma dos mapas.} Com a balanceada, a soma média do mapa de SNR é 14,0 (16,2 na atual; $-1{,}5$ na normal), ou seja, SNR tem o que explicar. A de BRD vai de 11,2 para 17,1, perto dos 18,6 da normal: o excesso de BRD nas baselines atuais, que “comia” parte da explicação de BRD, desaparece. \textbf{Região mais importante igual à do mapa atual:} 82\% com C, 72\% com D e 66\% com B.

\section{O que muda nas métricas}
Consistency ($k = 2$, valor observado) e complexidade (número efetivo de regiões, de 36), recalculadas com cada conjunto de mapas:
\begin{center}\small
\setlength{\tabcolsep}{4pt}
\begin{tabular}{lrrrrrrrrrrr}
\toprule
\textbf{Consistency} & \textbf{SNR} & \textbf{AF} & \textbf{IAVB} & \textbf{BRE} & \textbf{BRD} & \textbf{CAP} & \textbf{CVP} & \textbf{STD} & \textbf{STE} & \textbf{Global} & \textbf{Sem SNR} \\
\midrule
A. Atual, 50 & 0,21 & 0,42 & 0,45 & 0,25 & 0,64 & 0,14 & 0,26 & 0,51 & 0,03 & 0,42 & 0,45 \\
C. Atual, 200 & 0,25 & 0,46 & 0,43 & 0,33 & 0,64 & 0,11 & 0,35 & 0,52 & 0,09 & 0,44 & 0,46 \\
D. Balanceada, 200 & 0,26 & 0,46 & 0,42 & 0,34 & 0,72 & 0,16 & 0,38 & 0,39 & 0,08 & 0,46 & 0,48 \\
B. Normal, 200 & (0,58) & 0,47 & 0,59 & 0,35 & 0,76 & 0,16 & 0,50 & 0,45 & 0,03 & 0,54 & 0,54 \\
\midrule
\textbf{Complexidade} & & & & & & & & & & & \\
A. Atual, 50 & 26,9 & 22,8 & 23,3 & 25,8 & 22,4 & 25,7 & 24,7 & 26,9 & 25,5 & 24,2 & \\
C. Atual, 200 & 26,9 & 22,7 & 23,0 & 25,9 & 22,3 & 26,1 & 24,9 & 26,9 & 25,7 & 24,2 & \\
D. Balanceada, 200 & 26,2 & 22,9 & 23,0 & 26,2 & 22,2 & 25,4 & 25,0 & 26,9 & 25,5 & 24,1 & \\
B. Normal, 200 & (25,7) & 22,8 & 22,6 & 26,2 & 23,1 & 25,7 & 24,8 & 27,4 & 25,6 & 24,3 & \\
\bottomrule
\end{tabular}
\end{center}
\begin{itemize}
\item \textbf{Os 200 sorteios, sozinhos, mudam pouco a Consistency} (A $\to$ C: 0,42 $\to$ 0,44). Algumas classes sobem (BRE, CVP, STE), outras ficam iguais ou caem (IAVB, CAP). Isso responde a uma dúvida antiga: o relatório de 21/10 apontava que a Consistency mistura a variação real do modelo com o erro de estimativa dos 50 sorteios. Medido agora, o erro de estimativa pesa pouco; a Consistency baixa de várias classes reflete principalmente o modelo usar evidências diferentes em ECGs diferentes.
\item \textbf{A referência pesa mais que os sorteios na Consistency:} com 200 sorteios nos dois casos, a balanceada dá 0,46 e a normal, 0,54 (fora SNR, 0,48 e 0,54), contra 0,44 da atual. Com a normal, as explicações ficam mais parecidas entre ECGs da mesma classe.
\item \textbf{As pontas não mudam em nenhuma opção:} BRD é sempre a mais consistente, e CAP e STE, as menos. No meio, há trocas: STD cai com a balanceada (0,52 $\to$ 0,39), e CVP e IAVB sobem com a normal.
\item \textbf{A complexidade praticamente não depende da escolha} (24,1 a 24,3 no global; mesmas classes nas pontas).
\end{itemize}

\partbar{Parte V · Decisão}

\section{Opções}
\begin{center}\small
\setlength{\tabcolsep}{4pt}
\begin{tabularx}{\linewidth}{lXXXX}
\toprule
& \textbf{A. Manter} & \textbf{B. Normal, 200} & \textbf{C. Atual, 200} & \textbf{D. Balanceada, 200} \\
\midrule
Pergunta do mapa & Diferença em relação à média do treino (mistura desbalanceada) & Diferença em relação a um coração saudável & Igual à A & Diferença em relação a um ECG médio que não privilegia nenhuma classe \\ \addlinespace
Ruído (piso) & 0,355 & 0,167 & cerca de 0,18 & 0,183 \\ \addlinespace
SNR & Normal & Sem sentido & Normal & Normal \\ \addlinespace
Peso de BRD na referência & 14 de 32 & 0 & 14 de 32 & 4 de 36 \\ \addlinespace
Consistency (sem SNR) & 0,45 & 0,54 & 0,46 & 0,48 \\ \addlinespace
Defesa no texto & Padrão do SHAP; citar os testes & Mais natural clinicamente; tirar SNR ou tratá-la à parte & Reduz o ruído; referência continua desbalanceada & Reduz o ruído; referência neutra entre as classes; SNR mantida \\
\bottomrule
\end{tabularx}
\end{center}

\textbf{Custo de B, C e D (igual nas três).} Os mapas já foram gerados nos testes. É preciso refazer a Faithfulness Correlation e a Estimate (cerca de 30 min cada), a sensibilidade (cerca de 5,5 h) e o Random Logit (cerca de 4,5 h), que refazem o GradientSHAP e ficam 4 vezes mais lentos com 200 sorteios, além de revisar os textos. A Consistency e a complexidade já estão calculadas (Parte IV).

\textbf{Tratamento de SNR na opção B.} Duas saídas possíveis: (i) explicar SNR com outra referência e as demais classes com a normal, o que mistura duas perguntas na mesma tabela; ou (ii) tirar SNR das métricas de explicação e explicar no texto por quê (a classe é a própria referência). A comparação com o F1 passaria a ter 8 classes.

\begin{infobox}[note]{Leitura dos testes para a decisão}
\begin{itemize}
\item \textbf{As conclusões principais do TCC se mantêm em qualquer opção:} os mapas mudam pouco, a complexidade não muda e as pontas da Consistency são as mesmas.
\item \textbf{Os 200 sorteios valem pela precisão,} não por mudar resultados: o piso cai pela metade, o que fortalece a sensibilidade (hoje, piso 5 vezes maior que a métrica), o Random Logit e a leitura da Consistency.
\item \textbf{Entre as referências,} a balanceada (D) resolve o problema concreto da atual (BRD em excesso, BRE e STE ausentes) e mantém as 9 classes comparáveis com o F1; a normal (B) responde à pergunta mais clínica, ao custo de SNR.
\item \textbf{Sugestão para a discussão com o orientador:} opção D como principal, com a referência normal (B) citada como análise complementar, já que os resultados com ela estão neste documento.
\end{itemize}
\end{infobox}

\textbf{Para levar ao orientador:} qual referência usar (A, B, C ou D); se B, como tratar SNR; e se o número de sorteios passa a 200, com o custo de cerca de 10 h de Kaggle para refazer a sensibilidade e o Random Logit.

\vspace{0.5em}
{\small\sloppy \textbf{Arquivos}
\begin{itemize}
\item Primeiro teste: \texttt{src/baseline\_check.py} (commit \texttt{c3e5420}); resultados em \url{results/modeloB_basecheck/}, com a prévia das métricas da opção B em \texttt{previa\_normal200/}.
\item Segundo teste: \texttt{src/baseline\_balanced.py} (commit \texttt{4d07014}); resultados em \url{results/modeloB_basebal/}, com as métricas das opções C e D e as distâncias entre os mapas de 200 sorteios em \texttt{previa\_200/}.
\item Mapas de 200 sorteios (fora do git): \url{output/modeloB_normal200/}, \url{output/modeloB_atual200/}, \url{output/modeloB_balanceada200/}.
\end{itemize}}

\end{document}
